\documentclass[11pt,a4paper]{article}

% ------------------------------------------------------------------ %
% Packages
% ------------------------------------------------------------------ %
\usepackage[margin=2.2cm]{geometry}
\usepackage{amsmath,amssymb,mathtools}
\usepackage{booktabs}
\usepackage{longtable}
\usepackage{array}
\usepackage{xcolor}
\usepackage{listings}
\usepackage{siunitx}
\usepackage{microtype}
\usepackage{enumitem}
\usepackage{algorithm}
\usepackage{algpseudocode}
\usepackage[hidelinks]{hyperref}

% ------------------------------------------------------------------ %
% Listing style for parameter snippets
% ------------------------------------------------------------------ %
\lstdefinestyle{cfg}{
  basicstyle=\ttfamily\footnotesize,
  breaklines=true,
  columns=fullflexible,
  frame=single,
  rulecolor=\color{black!30},
  showstringspaces=false,
}
\lstset{style=cfg}

% Tighter algorithmic display
\algrenewcommand\algorithmicrequire{\textbf{Input:}}
\algrenewcommand\algorithmicensure{\textbf{Output:}}
\newcommand{\pyconst}[1]{\textcolor{black!60!blue}{\texttt{\footnotesize #1}}}
\newcommand{\dictkey}[1]{\textcolor{black!60!green}{\texttt{\footnotesize #1}}}

% ------------------------------------------------------------------ %
% Convenience macros
% ------------------------------------------------------------------ %
\newcommand{\Vm}{V_{m}}
\newcommand{\Vmsub}[1]{V_{m,#1}}
\newcommand{\Iion}{I_{\text{ion}}}
\newcommand{\Iionsub}[1]{I_{\text{ion},#1}}
\newcommand{\Iext}{I_{\text{ext}}}
\newcommand{\Dten}{\mathbf{D}}
\newcommand{\grad}{\nabla}
\newcommand{\divg}{\nabla\!\cdot}
\newcommand{\Aimp}{\mathsf{A}}
\newcommand{\Bimp}{\mathsf{B}}
\newcommand{\Mmass}{\mathsf{M}}
\newcommand{\Kstif}{\mathsf{K}}
\newcommand{\Lop}{\mathsf{L}}
\newcommand{\code}[1]{\texttt{\small #1}}

\title{Solver Specification:\\
       \texorpdfstring{\texttt{highOrderElectroActivationFoamImplicitJfNK}}{highOrderElectroActivationFoamImplicitJfNK}\\
       \large A high-order, implicit, Jacobian-free Newton--Krylov
       monodomain solver in OpenFOAM with an Eigen linear-algebra back-end}
\author{Pablo}
\date{\today}

\begin{document}
\maketitle

\begin{abstract}
This report specifies the algorithms implemented in the OpenFOAM solver
\code{highOrderElectroActivationFoamImplicitJfNK}. The solver
advances the monodomain equations of cardiac electrophysiology coupled
to the ten Tusscher--Noble--Noble--Panfilov (TNNP, 2004) ionic cell
model. The diffusion operator is discretised with the cell-centred
finite-volume framework of OpenFOAM, optionally upgraded to a Local
Radial-Extrapolation (LRE) high-order reconstruction; the ionic source
is evaluated either at cell centres or at LRE cell quadrature points, with the
reaction ODE integrated at the quadrature points, once per cell and
reconstructed, or by a front-aware hybrid that does the former only on the
propagating front and the latter elsewhere (Section~\ref{sec:stateint});
time integration follows a $\theta$-scheme (backward Euler or
Crank--Nicolson); the resulting nonlinear coupled system can be solved
by Picard fixed-point iteration, a diagonal $\Iion$ linearisation, or a
Jacobian-Free Newton--Krylov (JFNK) method with Armijo backtracking line
search, linear extrapolation initial guess and a physical-bracket clamp
on the input to the ionic ODE; the JFNK inner linear system is solved by
an in-house restarted GMRES with modified Gram--Schmidt orthogonalisation
built on top of the Eigen sparse linear algebra library; the per-cell
ionic ODE uses OpenFOAM's adaptive RKF45
integrator with configurable tolerances; activation times are tracked
by linear interpolation across the first threshold crossing and
benchmark points are sampled by inverse-distance weighting. All
algorithmic parameters are exposed through the
\code{constant/spatialIntegrationProperties} and
\code{constant/timeIntegrationProperties} dictionaries, and the
companion Python driver \code{run\_convergence.py} generates these
dictionaries from a single configuration file. This solver is
algorithmically identical to its PETSc-based sibling
\code{highOrderElectroActivationFoamImplicitJfNK\_PETSc}; the only
difference is the linear-algebra library (Eigen vs.\ PETSc).
\end{abstract}

\tableofcontents
\bigskip

% ================================================================== %
\section{Notation}
% ================================================================== %

\begin{center}
\begin{tabular}{@{}lll@{}}
\toprule
Symbol & Meaning & Units \\
\midrule
$\Omega \subset \mathbb{R}^{d}$ & Spatial domain ($d \in \{2,3\}$) & \si{\meter} \\
$\partial\Omega$ & Insulated boundary (zero-flux Neumann) & --- \\
$t \in (0, T]$ & Time & \si{\second} \\
$\Vm(\mathbf{x}, t)$ & Transmembrane voltage & \si{\volt} \\
$\mathbf{s}(\mathbf{x}, t) \in \mathbb{R}^{17}$ & TNNP ionic states & mixed \\
$\Iion(\Vm, \mathbf{s})$ & Total ionic current density & \si{\volt\per\second} \\
$\Iext(\mathbf{x}, t)$ & External stimulus current density & \si{\ampere\per\meter\cubed} \\
$\chi$ & Membrane surface-to-volume ratio & \si{\per\meter} \\
$C_{m}$ & Membrane capacitance per unit area & \si{\farad\per\meter\squared} \\
$\Dten(\mathbf{x})$ & Bulk conductivity tensor & \si{\siemens\per\meter} \\
$\theta \in \{0.5, 1\}$ & Time-stepping parameter (CN or BE) & --- \\
$\Delta t$ & Outer time-step length & \si{\second} \\
$N_{c}$ & Number of mesh cells & --- \\
$\mathsf{V}^{n} \in \mathbb{R}^{N_{c}}$ & Vector of cell-averaged $\Vm$ at time $t^{n}$ & \si{\volt} \\
$\Mmass, \Kstif, \Lop$ & Mass, stiffness, scaled diffusion matrix & --- \\
$\Aimp, \Bimp$ & Implicit LHS / RHS operators of the $\theta$-scheme & --- \\
$\mathsf{F}(\mathsf{x})$ & Discrete nonlinear residual to be zeroed & --- \\
\bottomrule
\end{tabular}
\end{center}

% ================================================================== %
\section{Governing equations}
% ================================================================== %

\subsection{Monodomain PDE}

The bulk voltage is propagated by the standard monodomain
reaction--diffusion equation~\cite{ttp2004,niederer2012}. In physical
current-density notation one may write
\begin{equation}
  \chi \, C_{m} \, \frac{\partial \Vm}{\partial t}
   \;=\;
   \divg\!\bigl( \Dten \grad \Vm \bigr)
   - \chi C_m\,\Iion(\Vm, \mathbf{s})
   + \Iext
   \quad \text{in } \Omega \times (0, T],
   \label{eq:monodomain}
\end{equation}
subject to insulated boundaries
$\mathbf{n} \cdot ( \Dten \grad \Vm ) = 0$ on
$\partial \Omega \times (0, T]$, and initial condition
$\Vm(\mathbf{x}, 0) = -\SI{84}{\milli\volt}$. The conductivity tensor
$\Dten$ is symmetric positive-definite and either uniform or read as a
\code{volTensorField} from disk. The factor $\chi C_m$ multiplying
$\Iion$ in~\eqref{eq:monodomain} reflects the convention used by the
ionic-model interface: the solver stores \code{Iion} as a voltage rate
with dimensions \si{\volt\per\second}, not as an
\si{\ampere\per\meter\cubed} volumetric current.

Dividing~\eqref{eq:monodomain} by $\chi C_{m}$ gives the form actually
used inside the solver,
\begin{equation}
  \frac{\partial \Vm}{\partial t}
   \;=\;
   \mathcal{L} \Vm
   \;+\;
   s(\Vm, \mathbf{s}),
   \qquad
   \mathcal{L} \;=\; \frac{1}{\chi C_{m}}\,\divg(\Dten \grad \cdot),
   \quad
   s \;=\; -\Iion + \frac{\Iext}{\chi C_{m}}.
   \label{eq:semi-discrete}
\end{equation}
This is exactly what \code{rhsFromIion} assembles in the solver.

\subsection{Ionic ODE}

At every spatial sampling point the TNNP state vector
$\mathbf{s} \in \mathbb{R}^{17}$ evolves according to
\begin{equation}
  \frac{\mathrm{d}\mathbf{s}}{\mathrm{d}t}
   \;=\;
   \mathbf{f}_{\text{TNNP}}\!\bigl(\Vm, \mathbf{s}; \texttt{tissue}\bigr),
   \qquad \mathbf{s}(0) = \mathbf{s}_{0},
   \label{eq:tnnp}
\end{equation}
where \code{tissue} selects \code{epicardialCells}, \code{mCells} or
\code{endocardialCells}. The 17 tracked states are
\[
  \mathbf{s} = (V,\ K_{i},\ Na_{i},\ Ca_{i},\ Xr_{1},\ Xr_{2},\ Xs,\ m,\ h,\ j,\ d,\ f,\ fCa,\ s,\ r,\ g,\ Ca_{SR}).
\]
The total ionic current returned to the PDE is the algebraic
\code{Iion\_cm} of TNNP~\cite{ttp2004}.

\subsection{Manufactured MMS equation used for convergence transfer}

The convergence driver in
\code{tutorials\_electro/highOrderManufacturedFDAImplicitJfNK} solves a
manufactured problem with the same numerical structure as the
electro-activation solver: implicit monodomain-like PDE, local state
ODEs, nonlinear ionic source, optional LRE reconstruction for $\Vm$ and
$\Iion$, and selectable \code{Picard}/\code{diagonalIion}/\code{JFNK}
nonlinear coupling. The analytical fields are
\begin{align}
  \Vm^{\star}(\mathbf{x},t) &= \sqrt{1+t}\,F(\mathbf{x}),\\
  u_1^{\star}(\mathbf{x},t) &= (1+t)\,G(\mathbf{x})+\Vm^{\star}(\mathbf{x},t),\\
  u_2^{\star}(\mathbf{x},t) &=
      \frac{1}{(1+t)\sqrt{G(\mathbf{x})}},\\
  u_3^{\star}(\mathbf{x},t) &= 0,
\end{align}
with, for example,
$F=\cos(\pi x)\cos(2\pi y)$ and $G=1+x\,y^2$ in 2D. In 3D,
$F=\cos(\pi x)\cos(2\pi y)\cos(3\pi z)$ and
$G=1+x\,y^2 z^3$. The state equations are
\begin{align}
  \dot u_1 &=
      (u_1+u_3-\Vm)^2 u_2^2
    + \frac{1}{2}(u_1+u_3-\Vm)u_2^2(\Vm-u_3),\\
  \dot u_2 &= -(u_1+u_3-\Vm)u_2^3,\\
  \dot u_3 &= 0.
\end{align}
The manufactured ionic source is
\begin{equation}
  \Iion^{\mathrm{MMS}}(\Vm,\mathbf{u})
  =
  -\frac{1}{2}(u_1+u_3-\Vm)u_2^2(\Vm-u_3)
  + \frac{\beta}{\chi C_m}(\Vm-u_3),
\end{equation}
and the PDE source used by the manufactured solver is
$-\Iion^{\mathrm{MMS}}$. The coefficient $\beta$ is the spatial
eigenvalue of the trigonometric manufactured field under the constant
conductivity tensor:
\[
  \beta =
  \begin{cases}
    -\pi^2 D_{xx}, & d=1,\\
    -\pi^2(D_{xx}+4D_{yy}), & d=2,\\
    -\pi^2(D_{xx}+4D_{yy}+9D_{zz}), & d=3.
  \end{cases}
\]
Thus the MMS case is not intended to reproduce TNNP physiology; it is a
controlled problem with a known exact solution that exercises the same
discrete operators and nonlinear solution paths used by this solver.

\subsection{External stimulus}

The external current is a time- and region-windowed step:
\begin{equation}
  \Iext(\mathbf{x}, t) \;=\;
  \begin{cases}
    I_{0} & \mathbf{x} \in \Omega_{\text{stim}}^{(b)},
            \;\; t \in [t_{0}^{(b)},\, t_{0}^{(b)} + \tau],
            \;\; b = 1, \dots, B, \\
    0     & \text{otherwise},
  \end{cases}
\end{equation}
with $I_{0} = \SI{50000}{\ampere\per\meter\cubed}$ and
$\tau = \SI{2}{\milli\second}$ as Niederer-benchmark defaults.

% ================================================================== %
\section{Time discretisation: the $\theta$-scheme}
% ================================================================== %

Time is advanced from $t^{n}$ to $t^{n+1} = t^{n} + \Delta t$ by
\begin{equation}
  \frac{\mathsf{V}^{n+1} - \mathsf{V}^{n}}{\Delta t}
   \;=\;
   \theta\, \Lop\, \mathsf{V}^{n+1}
   + (1-\theta)\, \Lop\, \mathsf{V}^{n}
   + \theta\, \mathbf{s}^{n+1}
   + (1-\theta)\, \mathbf{s}^{n}.
   \label{eq:theta}
\end{equation}

\begin{center}
\begin{tabular}{@{}lcl@{}}
\toprule
\code{implicitScheme} & $\theta$ & Scheme \\
\midrule
\code{backwardEuler} & $1$   & First-order, L-stable \\
\code{crankNicolson} & $1/2$ & Second-order, A-stable \\
\bottomrule
\end{tabular}
\end{center}

After multiplication by $\Delta t^{-1}$ the implementation forms
\begin{equation}
  \bigl[\Mmass/\Delta t \,-\, \theta\, \Lop\bigr]\, \mathsf{V}^{n+1}
   \;=\;
   \bigl[\Mmass/\Delta t \,+\, (1-\theta)\, \Lop\bigr]\, \mathsf{V}^{n}
   \;+\; \theta\, \mathbf{s}^{n+1}
   \;+\; (1-\theta)\, \mathbf{s}^{n},
   \label{eq:implicit-system}
\end{equation}
with the two Eigen \code{SparseMatrix<scalar,RowMajor>} operators
\begin{equation}
  \Aimp \;=\; \Mmass/\Delta t \,-\, \theta\, \Lop,
   \qquad
  \Bimp \;=\; \Mmass/\Delta t \,+\, (1-\theta)\, \Lop.
   \label{eq:AB}
\end{equation}
$\Aimp$ and $\Bimp$ are rebuilt at every time step by copying $\Mmass$,
scaling by $\Delta t^{-1}$ and adding $\pm \theta \Lop$. When
$\theta = 1$ (BE) the explicit half of the diffusion vanishes and
$\Bimp = \Mmass/\Delta t$. The scaled diffusion is
$\Lop = \Kstif / (\chi C_{m})$.

% ================================================================== %
\section{Spatial discretisation}
% ================================================================== %

\subsection{Standard orthogonal FV Laplacian}

\code{assembleStandardOrthogonalStiffnessMatrix} computes the classical
5/7-point Laplacian. Per internal face $f$ with face area $\mathbf{S}_{f}$,
\begin{equation}
  a_{f} \;=\;
   \frac{ \lvert \mathbf{S}_{f} \rvert\,
          \bigl( \mathbf{n}_{f} \cdot \Dten_{f} \cdot \hat{\mathbf{d}}_{f} \bigr) }
        { \lvert \mathbf{d}_{f} \rvert },
   \quad
   \Dten_{f} = \tfrac{1}{2}\,(\Dten_{\text{own}} + \Dten_{\text{nei}}),
\end{equation}
contributing $-a_{f}/V_{\text{own}}$ to $\Kstif[\text{own},\text{own}]$,
$+a_{f}/V_{\text{own}}$ to $\Kstif[\text{own},\text{nei}]$, and the
transposed pair to the neighbour row. \code{zeroGradient}/\code{empty}
patches contribute zero flux. \code{fixedValue} patches contribute to
the owner diagonal only.

\subsection{High-order LRE diffusion}

When \code{useHighOrder\_Vm = true} and $d>1$,
\code{assembleHighOrderStiffnessMatrix} integrates the flux on a face
quadrature with reconstructed gradients:
\begin{equation}
  \Phi_{f} \;=\;
   \sum_{q \in \mathcal{Q}_{f}}
       w_{q}\, \lvert \mathbf{S}_{f} \rvert\,
       \mathbf{n}_{f} \cdot \Dten_{f} \cdot
       \Bigl(
         \sum_{j \in \mathcal{S}_{f}}
            \mathbf{g}_{f,q,j}\, \mathsf{V}_{j}
       \Bigr),
   \label{eq:HO-flux}
\end{equation}
where $\mathcal{S}_{f}$ is the LRE face stencil ($\sim 15$--$50$ cells)
and $\mathbf{g}_{f,q,j} = $~\code{QRGradFaceGPCoeffs}. The
contribution of each $\mathsf{V}_{j}$ is added to row \code{own} with
weight $+\Phi_{f,j}/V_{\text{own}}$ and to row \code{nei} with weight
$-\Phi_{f,j}/V_{\text{nei}}$, yielding a negative semi-definite
Laplacian.

\subsection{Mass matrices}

Two options:
\paragraph{Lumped (\code{massMatrix lumped}).}
\code{assembleDiagonalMassMatrix} writes $\Mmass = I$.

\paragraph{Consistent HO (\code{massMatrix consistent}).}
\code{assembleConsistentMassMatrixHO} fills
\begin{equation}
  \Mmass[i, j] \;=\;
   \sum_{q \in \mathcal{Q}_{i}}
       \frac{w_{q}}{\sum_{q'} w_{q'}}\,
       \phi_{j}(\mathbf{x}_{q}),
   \quad
   \phi_{j}(\mathbf{x}_{q})
    = \delta_{ij}
    + \mathbf{g}_{i,j} \cdot \mathbf{d}_{q}
    + \tfrac{1}{2}\, \mathbf{H}_{i,j} \!:\! (\mathbf{d}_{q} \otimes \mathbf{d}_{q})
    + \tfrac{1}{6}\, \mathbf{T}^{(3)}_{i,j} \!\vdots\!
        (\mathbf{d}_{q} \otimes \mathbf{d}_{q} \otimes \mathbf{d}_{q}),
   \label{eq:HO-mass}
\end{equation}
with $\mathbf{d}_{q} = \mathbf{x}_{q} - \mathbf{C}_{i}$. A safety guard
silently downgrades to the lumped form when
\code{useHighOrder\_Vm = false} or $d \le 1$.

\subsection{High-order reconstruction at ionic quadrature points}

If \code{useHighOrder\_Iion = true}, the ionic ODE is integrated at
cell quadrature points. Vm is reconstructed there by the same Taylor
expansion as in~\eqref{eq:HO-mass}:
\begin{equation}
  \Vm(\mathbf{x}_{q})
   \;\approx\;
   \Vmsub{c} + (\grad \Vm)_{c} \cdot \mathbf{d}_{q}
   + \tfrac{1}{2}\, \mathbf{H}_{c} \!:\! (\mathbf{d}_{q} \otimes \mathbf{d}_{q})
   + \tfrac{1}{6}\, \mathbf{T}^{(3)}_{c} \!\vdots\!
       (\mathbf{d}_{q} \otimes \mathbf{d}_{q} \otimes \mathbf{d}_{q}),
   \label{eq:HO-vm}
\end{equation}
implemented in \code{reconstructVmAtIionIntegrationPoints}.

\subsection{High-order cell averaging of $\Iion$}

\code{averageIionIntegrationPointsToCells} contracts the per-QP $\Iion$
back to cells:
\begin{equation}
  \Iionsub{i} \;=\;
   \frac{\sum_{q \in \mathcal{Q}_{i}} w_{q}\, \Iionsub{q}}
        {\sum_{q \in \mathcal{Q}_{i}} w_{q}}.
\end{equation}

% ================================================================== %
\section{State integration at the ionic quadrature points}
\label{sec:stateint}
% ================================================================== %

The high-order quadrature~\eqref{eq:HO-vm} evaluates $\Iion(\Vm,\mathbf{s})$
at the cell quadrature points $\mathbf{x}_q$, so it needs the ionic state
vector $\mathbf{s}$ there as well. \emph{How} the states reach the quadrature
points is governed by \dictkey{stateIntegrationMode} and, orthogonally, by
\dictkey{frontAwareHybrid}. This choice turns out to be decisive for a
propagating cardiac front, and this section documents both the algorithm and
the reason it matters.

% ------------------------------------------------------------------ %
\subsection{Two modes}
% ------------------------------------------------------------------ %

Let $\mathcal{Q}_i=\{\mathbf{x}_q\}$ be the quadrature set of cell $\Omega_i$
with weights $w_q$, and let $\mathbf{s}_c$ denote a cell-centred state.

\paragraph{\code{gaussPointODE} (legacy).}
The stiff ODE~\eqref{eq:tnnp} is integrated \emph{independently at every
quadrature point}, driven by the locally reconstructed potential:
\begin{equation}
  \mathbf{s}_q^{\,n+1}
   = \mathrm{ODEsolve}\!\bigl(\mathbf{s}_q^{\,n},\,\Vm(\mathbf{x}_q);\Delta t\bigr),
  \qquad
  \Iionsub{q}=\phi\bigl(\Vm(\mathbf{x}_q),\mathbf{s}_q^{\,n+1}\bigr).
  \label{eq:gpode}
\end{equation}
Cost: $N_q$ stiff integrations per cell.

\paragraph{\code{cellCentredReconstruct}.}
The ODE is integrated \emph{once per cell} at the cell centre, and the
resulting cell-centred states are reconstructed to the quadrature points by a
dedicated LRE (\code{LREInterp\_states}, order set by
\dictkey{highOrderCoeffs/LRECoeffs\_states}):
\begin{equation}
  \mathbf{s}_c^{\,n+1}
   = \mathrm{ODEsolve}\!\bigl(\mathbf{s}_c^{\,n},\,\overline{V}_{m,c};\Delta t\bigr),
  \quad
  \widehat{\mathbf{s}}(\mathbf{x}_q)
   = \mathcal{R}^{(p)}\!\bigl[\{\mathbf{s}_{c'}^{\,n+1}\}\bigr](\mathbf{x}_q),
  \quad
  \Iionsub{q}=\phi\bigl(\Vm(\mathbf{x}_q),\widehat{\mathbf{s}}(\mathbf{x}_q)\bigr).
  \label{eq:ccrec}
\end{equation}
Cost: one stiff integration per cell -- a factor $N_q$ cheaper, which for the
expensive TNNP integrations is the motivation. The two modes differ only in the
argument of the non-linear map $\Vm\!\mapsto\!\mathbf{s}$: pointwise
$\Vm(\mathbf{x}_q)$ in~\eqref{eq:gpode}, versus the cell average
$\overline{V}_{m,c}$ in~\eqref{eq:ccrec}.

% ------------------------------------------------------------------ %
\subsection{Why the reaction ODE must be solved at the Gauss points on the front}
\label{sec:jensen}
% ------------------------------------------------------------------ %

For a smooth reaction field the two modes agree to the reconstruction order,
and \code{cellCentredReconstruct} is preferable for its speed. On a
\emph{propagating cardiac front} they do not agree, and only the pointwise
evaluation~\eqref{eq:gpode} is admissible. The reason is the strong
non-linearity of the sodium activation.

The upstroke is carried by
$I_{\mathrm{Na}}=g_{\mathrm{Na}}\,m^{3}h\,j\,(\Vm-E_{\mathrm{Na}})$, whose
activation gate relaxes towards the steep sigmoid
\begin{equation}
  m_\infty(\Vm)=\Bigl[1+\exp\!\bigl((V_{1/2}-\Vm)/k\bigr)\Bigr]^{-1},
  \qquad V_{1/2}\approx\SI{-40}{mV},\; k\approx\SI{6}{mV}.
  \label{eq:minf}
\end{equation}
Below its inflection ($\Vm<V_{1/2}$) -- exactly the sub-threshold foot where
ignition and conduction are decided -- $m_\infty$ is \emph{convex}. Writing
$\langle\cdot\rangle_i=\bigl(\sum_q w_q\,\cdot\bigr)/\!\sum_q w_q$ for the
cell-quadrature average, so that $\overline{V}_{m,i}=\langle\Vm\rangle_i$, the
two modes evaluate the gate as
\begin{equation}
  \underbrace{\bigl\langle m_\infty(\Vm)\bigr\rangle_i}_{\text{\code{gaussPointODE}}}
  \;\;\ge\;\;
  \underbrace{m_\infty\!\bigl(\langle\Vm\rangle_i\bigr)}_{\text{\code{cellCentredReconstruct}}},
  \label{eq:jensen}
\end{equation}
by Jensen's inequality for the convex $m_\infty$. Since
$I_{\mathrm{Na}}\propto m^{3}$, the gap is amplified cubically: the cell-average
path \emph{systematically under-estimates} the regenerative inward current
precisely where $\Vm$ varies within a cell, i.e.\ at the wavefront and the
stimulus edge. Reconstructing the states after the ODE
(as in~\eqref{eq:ccrec}) cannot repair this -- the bias is committed
\emph{inside} the ODE, whose input was already the averaged potential.

\paragraph{A quantitative estimate.}
Let $\Delta V=\max_q\Vm(\mathbf{x}_q)-\min_q\Vm(\mathbf{x}_q)$ be the sub-cell
range of $\Vm$ across a cell straddling the foot of the upstroke, and treat
$\Vm$ as roughly uniform over the cell (variance $\sigma^2\simeq\Delta V^2/12$).
A midpoint expansion of the average gives
\begin{equation}
  \bigl\langle m_\infty\bigr\rangle
  - m_\infty(\overline V)
  \;\simeq\; \tfrac12\,m_\infty''(\overline V)\,\sigma^2
  \;=\; m_\infty''(\overline V)\,\frac{\Delta V^{2}}{24}.
\end{equation}
On the exponential foot $m_\infty\!\approx\!e^{(\Vm-V_{1/2})/k}$, so
$m_\infty''/m_\infty=1/k^{2}$ and, using $I_{\mathrm{Na}}\propto m^{3}$,
\begin{equation}
  \boxed{\;\frac{\delta I_{\mathrm{Na}}}{I_{\mathrm{Na}}}
  \;\approx\; 3\cdot\frac{\Delta V^{2}}{24\,k^{2}}
  \;=\; \frac{\Delta V^{2}}{8\,k^{2}}\; }.
  \label{eq:gap}
\end{equation}
With $k=\SI{6}{mV}$ and a sub-cell range $\Delta V\sim\SI{30}{mV}$ across an
under-resolved front, \eqref{eq:gap} gives
$\delta I_{\mathrm{Na}}/I_{\mathrm{Na}}\approx 900/288\approx 3$ -- the linearised
estimate is already $O(1)$, i.e.\ the cell-average evaluation under-shoots the
sodium current by a factor of order unity or more. Two failure modes follow,
both observed and documented in the companion note
\code{stateIntegrationMode\_conductionDiagnosis.tex}:
\begin{enumerate}[topsep=2pt,itemsep=1pt]
  \item \textbf{Ignition failure:} a nominal stimulus is rendered effectively
    sub-threshold; $\Vm$ rises but the regenerative upstroke never fires.
  \item \textbf{Conduction block:} once ignited, each front cell supplies too
    little inward charge to bring its downstream neighbours to threshold -- a
    numerically induced source--sink mismatch that stalls the wave.
\end{enumerate}
Refining $\Delta t$ makes the front sharper (larger $\Delta V$) and thus
\emph{aggravates} the deficit, which is why the pure \code{cellCentredReconstruct}
path fails at small $\Delta t$ while a coarse-$\Delta t$ over-shoot can mask it.
The conclusion is unavoidable: at high order on a sharp front the reaction ODE
\emph{must} be solved at the Gauss points of the front cells; only the remaining
(smooth) volumes may be reconstructed cheaply.

% ------------------------------------------------------------------ %
\subsection{Front-aware hybrid}
\label{sec:hybrid}
% ------------------------------------------------------------------ %

\dictkey{frontAwareHybrid} (default \code{false}, orthogonal to
\dictkey{stateIntegrationMode}) realises exactly that conclusion: it solves the
ODE at the Gauss points in the thin band of \emph{front} cells and at the cell
centre (with reconstruction) everywhere else. A cell $i$ is a \emph{front} cell
when its sub-cell $\Vm$ range exceeds \dictkey{hybridFrontVmRange} or it is
stimulated,
\begin{equation}
  \text{front}_i \;\Longleftrightarrow\;
  \Bigl(\max_q\Vm(\mathbf{x}_q)-\min_q\Vm(\mathbf{x}_q)\Bigr)
     > \dictkey{hybridFrontVmRange}
  \;\;\vee\;\; (\Iext)_i\neq 0 .
  \label{eq:frontflag}
\end{equation}
The persistent states are carried at the Gauss points (as in
\code{gaussPointODE}); the smooth-cell reconstruction derives its cell values
by quadrature-averaging that carrier. Per nonlinear evaluation the branch is:
\begin{enumerate}[topsep=2pt,itemsep=1pt,label=(\alph*)]
  \item front cells: one ODE per Gauss point, local $\Vm$ (as~\eqref{eq:gpode});
  \item smooth cells: one ODE at the cell centre from the cell-average state;
  \item reconstruct the smooth-cell states to their Gauss points, front Gauss
        states left untouched;
  \item evaluate $\Iion$ at all Gauss points and average to the cell.
\end{enumerate}

\paragraph{Freezing per time step (avoiding a limit cycle).}
The front classification~\eqref{eq:frontflag} is a \emph{solution-dependent
discrete switch}. If recomputed inside the nonlinear loop, cells near the
threshold flip front$\leftrightarrow$smooth between iterations; because the two
branches compute $\Iion$ by different paths, the source term jumps discretely
and the Picard fixed-point map becomes discontinuous, producing a non-convergent
\emph{period-2 limit cycle} (observed as $\text{iter}_{k+2}\equiv\text{iter}_k$
in \code{nonlinearResiduals.dat}). The classification is therefore computed
\emph{once per time step} from the committed $\Vm^{n}$ and \emph{frozen} for the
whole nonlinear solve, and the band is dilated by \dictkey{hybridFrontDilation}
cell layers (via \code{mesh.cellCells()}) so the front stays inside the
Gauss-resolved zone as it advances. Freezing is safe because the front moves far
less than one cell per step (for the 2D benchmark, $\sim\!\SI{0.09}{mm/ms}$, i.e.
$\sim\!0.009$ cell per step at $\Delta t=\SI{e-4}{s}$, $\Delta x=\SI{0.5}{mm}$).

% ------------------------------------------------------------------ %
\subsection{State reconstruction limiter}
\label{sec:limiter}
% ------------------------------------------------------------------ %

High-order ($p3$) reconstruction of the stiff TNNP states across an
under-resolved front overshoots (Gibbs), and a state driven far outside its
physical range can overflow the ionic algebra (observed as a PETSc
\code{SIGFPE}). \dictkey{stateReconstructLimiter} (default \code{true}, active
for \code{cellCentredReconstruct} and the hybrid smooth cells, independent of
\dictkey{frontAwareHybrid}) applies a Barth--Jespersen-style bound: each
reconstructed value is clipped to the range of the cell and its face neighbours,
\begin{equation}
  \widehat{s}(\mathbf{x}_q)
  \;\leftarrow\;
  \mathrm{clip}\!\Bigl(\widehat{s}(\mathbf{x}_q),\;
     \min_{c'\in\{i\}\cup\mathcal{N}(i)} s_{c'},\;
     \max_{c'\in\{i\}\cup\mathcal{N}(i)} s_{c'}\Bigr).
  \label{eq:limiter}
\end{equation}
No new extremum is created, so the reconstructed states stay within the
physically realised (ODE-evolved) range and the ionic algebra never sees a
pathological input; in smooth regions the reconstruction already lies within the
range and the limiter is inactive, preserving the order there.

% ------------------------------------------------------------------ %
\subsection{Dictionary keys}
% ------------------------------------------------------------------ %

\begin{center}
\begin{tabular}{@{}lll@{}}
\toprule
Key (in \code{implicitHOCoeffs}) & Meaning & Default \\
\midrule
\code{stateIntegrationMode} & \code{gaussPointODE} / \code{cellCentredReconstruct} & \code{cellCentredReconstruct} \\
\code{frontAwareHybrid} & Gauss ODEs on the front, reconstruct elsewhere & \code{false} \\
\code{hybridFrontVmRange} & sub-cell $\Vm$ range flagging a front cell [V] & $5\times10^{-3}$ \\
\code{hybridFrontDilation} & cell layers added around the front band & $2$ \\
\code{stateReconstructLimiter} & Barth--Jespersen limiter on the reconstruction & \code{true} \\
\bottomrule
\end{tabular}
\end{center}
\code{LRECoeffs\_states} (in \code{highOrderCoeffs}) sets the reconstruction
order; it is required by \code{cellCentredReconstruct} and by the hybrid.

% ================================================================== %
\section{Nonlinear solver}
% ================================================================== %

At each outer time step the unknown is the new voltage vector
$\mathsf{x}=\mathsf{V}^{n+1}$. For any candidate $\mathsf{x}$ the solver
first reconstructs the corresponding $\Vm$ field, then calls
\code{evaluateNonlinearFields}. That routine resets the ionic states to
their time-step-start values $\mathbf{s}^{n}$, integrates the local ODEs
from $t^n$ to $t^{n+1}$ using the candidate voltage, obtains
$\Iion(\mathsf{x})$, and optionally averages quadrature-point values
back to cells. This gives the nonlinear residual
\begin{equation}
  \mathsf{F}(\mathsf{x})
   \;=\;
   \Aimp\, \mathsf{x}
   \;-\;
   \mathsf{r}\!\bigl( \Iion(\mathsf{x}) \bigr),
   \qquad
   \mathsf{x} = \mathsf{V}^{n+1},
   \label{eq:Fdef}
\end{equation}
with
\begin{equation}
  \mathsf{r}(\Iion)
  =
  \Bimp\,\mathsf{V}^{n}
  + \theta\,\Bigl[-\Iion(\mathsf{x})+\Iext^{n+1}/(\chi C_m)\Bigr]
  + (1-\theta)\,\Bigl[-\Iion^{n}+\Iext^{n}/(\chi C_m)\Bigr].
  \label{eq:rhsIion}
\end{equation}
The three nonlinear methods differ only in how they drive
$\mathsf{F}(\mathsf{x})$ toward zero.

\subsection{Shared residuals and stopping tests}

Every method records the same diagnostics:
\begin{align}
  R_{\mathrm{cpl}} &=
    \frac{\|\mathsf{F}(\mathsf{x}_{k})\|_2}
         {\|\mathsf{x}_{k}\|_2+\epsilon},\\
  R_{\Vm} &=
    \frac{\|\mathsf{V}_{k}-\mathsf{V}_{k-1}\|_2}
         {\|\mathsf{V}_{k-1}\|_2+\epsilon},\\
  R_{\Iion} &=
    \frac{\|I_{\mathrm{ion},k}-I_{\mathrm{ion},k-1}\|_2}
         {\|I_{\mathrm{ion},k-1}\|_2+\epsilon},\\
  R_{\mathbf{s}} &=
    \max_j
    \frac{\|\mathbf{s}_{j,k}-\mathbf{s}_{j,k-1}\|_2}
         {\|\mathbf{s}_{j,k-1}\|_2+\epsilon}.
\end{align}
The actual implementation uses \code{relativeL2Difference} for
$R_{\Vm}$ and $R_{\Iion}$, and \code{maxStateRelativeL2Difference} for
the state maximum. A nonlinear iteration can stop only after
\code{minNonlinearIterations}; convergence then requires
\code{nonlinearTolerance}, \code{nonlinearVmTolerance} and
\code{nonlinearIionTolerance}. The state criterion
\code{nonlinearStatesTolerance} is included only when
\code{nonlinearRequireStatesConvergence=true}.

\subsection{Picard fixed-point method}

The Picard branch treats the source as lagged during the PDE solve. At
iteration $k$ the states and ionic source are computed from
$\mathsf{x}_k$, then the linear PDE
\begin{equation}
  \Aimp\,\mathsf{V}_{*}
  =
  \mathsf{r}\!\bigl(\Iion(\mathsf{x}_{k})\bigr)
\end{equation}
is solved with the Eigen sparse solver selected by \code{linearSolver}.
The next candidate is relaxed,
\begin{equation}
  \mathsf{x}_{k+1}
  =
  \mathsf{x}_{k}
  + \omega(\mathsf{V}_{*}-\mathsf{x}_{k}),
  \qquad
  \omega=\code{nonlinearRelaxation}.
\end{equation}
After the update the solver re-evaluates states and $\Iion$ at
$\mathsf{x}_{k+1}$ and checks the residuals above. In this
electro-activation solver, unlike a pure MMS Picard study, a failed
Picard solve rolls back unless \code{nonlinearAcceptUnconverged=true}.

\paragraph{Advantages.}
Picard is simple, cheap per nonlinear iteration, and uses the same
linear PDE solver as the classical implicit formulation. It is a useful
baseline and is often adequate when $\Delta t$ is small or the ionic
source is weakly nonlinear over one time step.

\paragraph{Disadvantages.}
The source derivative with respect to $\Vm$ is ignored in the PDE
matrix, so convergence can be slow near the action-potential upstroke or
for large $\Delta t$. A small residual may require many iterations, and
under-relaxation may be needed for robustness.

\subsection{Diagonal \texorpdfstring{$\Iion$}{Iion} linearisation}

The \code{diagonalIion} branch is Picard plus a local finite-difference
approximation of the source derivative. Define
\[
  S(\Vm) = -\Iion(\Vm,\mathbf{s}(\Vm)) + \Iext/(\chi C_m).
\]
The solver approximates, cell by cell,
\begin{equation}
  S'_i(V_{m,k})
  \;\approx\;
  -\frac{
      I_{\mathrm{ion},i}(V_{m,k}+\varepsilon_d)
      - I_{\mathrm{ion},i}(V_{m,k})
    }{\varepsilon_d},
  \qquad
  \varepsilon_d=\code{diagonalIionEpsilon}.
\end{equation}
The linearised residual
$\Aimp\mathsf{x}-\theta S(\mathsf{x})$ gives the system
\begin{equation}
  \left(\Aimp-\theta\,\mathrm{diag}(S'_k)\right)\mathsf{x}_{k+1}
  =
  \mathsf{r}\!\bigl(\Iion(\mathsf{x}_{k})\bigr)
  - \theta\,\mathrm{diag}(S'_k)\,\mathsf{x}_{k}.
  \label{eq:diagonalIion}
\end{equation}
The update can still be relaxed with \code{nonlinearRelaxation}, and
the same convergence and rollback logic as Picard is used.

\paragraph{Advantages.}
It captures the dominant local stiffness of $\Iion(\Vm)$ without
building a full Jacobian. It often reduces the number of nonlinear
iterations relative to Picard while preserving a sparse scalar PDE
matrix.

\paragraph{Disadvantages.}
Only the cell-local diagonal sensitivity is included. Couplings through
ODE states, high-order quadrature averaging and spatial reconstruction
are not represented in the matrix. The finite-difference step
\code{diagonalIionEpsilon} must be chosen above ODE/noise levels but
small enough to be a useful derivative.

\subsection{Jacobian-Free Newton--Krylov}

JFNK applies Newton to~\eqref{eq:Fdef}:
\[
  \mathsf{J}(\mathsf{x}_k)\boldsymbol{\delta}_k
  =
  -\mathsf{F}(\mathsf{x}_k),
  \qquad
  \mathsf{x}_{k+1}=\mathsf{x}_k+\alpha_k\boldsymbol{\delta}_k,
\]
but never assembles $\mathsf{J}$. GMRES sees only matrix-vector
products approximated by
\begin{equation}
  \mathsf{J}(\mathsf{x}_{k})\, \mathbf{v}
   \;\approx\;
   \frac{\mathsf{F}(\mathsf{x}_{k} + \varepsilon\, \mathbf{v})
         - \mathsf{F}(\mathsf{x}_{k})}{\varepsilon},
   \qquad
   \varepsilon \;=\; \varepsilon_{0}\,
                     \frac{1 + \lVert \mathsf{x}_{k} \rVert_{2}}
                          {\lVert \mathbf{v} \rVert_{2}},
   \label{eq:fdjac}
\end{equation}
with $\varepsilon_{0}=\code{jfnkEpsilon}$. Every evaluation of
$\mathsf{F}$ re-runs the ODEs and recomputes $\Iion$ from the same
time-step-start state snapshot. This is expensive but gives a genuinely
coupled nonlinear residual.

\subsection{In-house restarted GMRES}

\code{solveGMRES} implements GMRES(m) with modified Gram--Schmidt. It
solves the Krylov least-squares problem for the Newton correction using
at most \code{jfnkMaxKrylovIterations} basis vectors per restart and at
most \code{jfnkMaxRestarts} restarts. The stopping test is controlled by
\code{jfnkLinearTolerance}. The method is ``matrix-free'' here because
the linear operator is the finite-difference shell in~\eqref{eq:fdjac},
not a stored sparse matrix.

\paragraph{Advantages.}
JFNK sees the full nonlinear dependence
$\Vm\mapsto\mathbf{s}(\Vm)\mapsto\Iion(\Vm,\mathbf{s})$ and usually has
far better nonlinear convergence than Picard in stiff regimes. It also
avoids explicitly forming a dense or complicated Jacobian.

\paragraph{Disadvantages.}
Each Krylov vector requires an additional residual evaluation, hence a
new ODE integration at every cell or quadrature point. The method is
sensitive to the finite-difference perturbation, ODE tolerances and
line-search choices. In this Eigen implementation the JFNK inner GMRES
has no preconditioner.

\subsection{Armijo backtracking line search}

The raw Newton update may increase the nonlinear residual when the
linear model is poor. If \code{jfnkLineSearch=true}, the solver tries
\[
  \mathsf{x}_{k+1}=\mathsf{x}_{k}+\alpha\boldsymbol{\delta}_{k},
  \qquad
  \alpha=\code{nonlinearRelaxation},\;
  \alpha/2,\;\alpha/4,\dots
\]
until the Armijo sufficient-decrease condition holds:
\begin{equation}
  \lVert \mathsf{F}(\mathsf{x}_{k}+\alpha\boldsymbol{\delta}_{k})
  \rVert_2^2
  \le
  (1-2c\alpha)\,\lVert \mathsf{F}(\mathsf{x}_{k})\rVert_2^2,
  \qquad c=\code{jfnkArmijoC}.
  \label{eq:armijo}
\end{equation}
Backtracking stops after \code{jfnkLineSearchMaxIter} attempts or when
the step length reaches \code{jfnkLineSearchAlphaMin}. This globalises Newton:
large Newton steps are accepted when they reduce the residual, otherwise
they are damped.

\subsection{Initial guess, ODE clamp and rollback}

The first guess for the time step is either $\mathsf{V}^{n}$ or the
linear extrapolation
\begin{equation}
  \mathsf{x}_{0}
  =
  \mathsf{V}^{n}
  + \frac{\Delta t}{\Delta t_{\mathrm{prev}}}
    \left(\mathsf{V}^{n}-\mathsf{V}^{n-1}\right),
  \label{eq:extrap}
\end{equation}
selected by \code{jfnkInitGuessOrder}. It is clamped to
$[\code{jfnkInitGuessVmMin},\code{jfnkInitGuessVmMax}]$. The same
bracket can also be applied to voltage values passed into the ODE solver
through \code{jfnkClampODEInput}; this protects the ionic model during
JFNK finite-difference probes.

If the nonlinear loop reaches its iteration budget without satisfying
the convergence tests, \code{nonlinearAcceptUnconverged} decides whether
the last iterate is committed. With the default \code{false}, the solver
rolls back $\Vm$, $\Iion$ and all ionic states to the beginning of the
time step. Rollback prevents a bad nonlinear solve from contaminating
the remaining activation simulation.

% ================================================================== %
\section{Linear solver: Eigen back-end}
% ================================================================== %

The wrapper \code{solveSparseSystem} drives either of two Eigen sparse
solvers on the explicit operator $\Aimp$ (Picard / diagonal paths):

\begin{center}
\begin{tabular}{@{}lll@{}}
\toprule
\code{linearSolver} key & Eigen type & Type \\
\midrule
\code{SparseLU}  & \code{Eigen::SparseLU} & Direct LU \\
\code{BiCGSTAB}  & \code{Eigen::BiCGSTAB} + \code{IncompleteLUT} & Iterative \\
\bottomrule
\end{tabular}
\end{center}

Both share the \code{tolerance} and \code{maxIterations} dictionary
keys (the former is ignored by \code{SparseLU} since it is direct).
The matrix-free JFNK branch never instantiates these wrappers; instead
it calls the in-house GMRES (Section~5.2) on the FD Jacobian shell.

% ================================================================== %
\section{Ionic ODE integration}
% ================================================================== %

For each cell or cell QP, the TNNP ODE~\eqref{eq:tnnp} is advanced by
the OpenFOAM adaptive integrator selected via the \code{solver} key.
Default \code{RKF45}~\cite{hairer2}. Keys read from
\code{constant/timeIntegrationProperties}:

\begin{center}
\begin{tabular}{@{}lll@{}}
\toprule
Key & Meaning & Default \\
\midrule
\code{solver} & Concrete ODE integrator & \code{RKF45} \\
\code{absTol} & Absolute tolerance per state & $10^{-9}$ \\
\code{relTol} & Relative tolerance per state & $10^{-6}$ \\
\code{maxSteps} & Cap on sub-steps per outer call & $10^{4}$ \\
\code{initialODEStep} & Initial sub-step (ms) & $10^{-5}$ \\
\bottomrule
\end{tabular}
\end{center}

\code{relTol} drives the noise floor of the matrix-free Jacobian:
$\eta_{\mathsf{F}} \sim $~\code{relTol}. Tightening \code{relTol} buys
a lower attainable outer tolerance at the cost of more RKF45 sub-steps
per Newton iteration.

% ================================================================== %
\section{TNNP numerical protection}
% ================================================================== %

The master switch \code{TNNPNumericalProtection} (default \code{false})
gates a local guard invoked at every external entry into the embedded TNNP
kernel (\code{calculateCurrent}, \code{solveODE\_start},
\code{solveODE\_end}, \code{derivatives}).

\begin{center}
\begin{tabular}{@{}lll@{}}
\toprule
Key & Type & Default \\
\midrule
\code{TNNPNumericalProtection} & Switch & \code{false} \\
\code{TNNPClampGates} & Switch & \code{true} \\
\code{TNNPClampVmForRates} & Switch & \code{true} \\
\code{TNNPLogNumericalProtection} & Switch & \code{true} \\
\code{TNNPConcentrationFloor} & scalar (mM) & $10^{-12}$ \\
\code{TNNPVMinForRates} & scalar (mV) & $-200$ \\
\code{TNNPVMaxForRates} & scalar (mV) & $+100$ \\
\code{TNNPVoltageZeroEps} & scalar (mV) & $10^{-8}$ \\
\code{TNNPMaxProtectionLogEntries} & label & $10^{6}$ \\
\code{TNNPProtectionLog} & word & \code{TNNP\_numericalProtection\_summary.dat} \\
\bottomrule
\end{tabular}
\end{center}

When protection is enabled, the guards applied in order are:
\begin{enumerate}[topsep=0pt,itemsep=2pt]
  \item \textbf{Non-finite $V$}: \code{!std::isfinite(V)} $\rightarrow$
        reset to $-86.2$ mV.
  \item \textbf{$V$ clamp for rates} (if \code{TNNPClampVmForRates}):
        clamp to $[\,$\code{TNNPVMinForRates}, \code{TNNPVMaxForRates}$\,]$.
  \item \textbf{$V=0$ singularity avoidance}: push $V$ off zero by
        \code{TNNPVoltageZeroEps} (sign preserving) when
        $\lvert V \rvert$ falls below that threshold.
  \item \textbf{Concentration floor} on $K_{i}, Na_{i}, Ca_{i}, Ca_{SR}$:
        non-finite or below-floor entries are reset to
        \code{TNNPConcentrationFloor}.
  \item \textbf{Gate clamp} on $Xr_{1}, Xr_{2}, Xs, m, h, j, d, f, fCa,
        s, r, g$ (if \code{TNNPClampGates}): non-finite to zero,
        clamp to $[0, 1]$.
\end{enumerate}

The embedded implementation stores a capped correction count and emits a
compact summary file at exit; it does not depend on external ionic-model
logging utilities.

% ================================================================== %
\section{Activation tracking and benchmarks}
% ================================================================== %

\paragraph{Activation time per cell.}
First-crossing of $V_{\text{thr}}$ in a time step is recorded by linear
interpolation:
\begin{equation}
  w \;=\; \mathrm{clip}\!\Bigl(
      \frac{V_{\text{thr}} - V^{n}}{V^{n+1} - V^{n}},
      \; 0, 1
    \Bigr),
   \qquad
   t_{\text{act}} \;=\; t^{n} + w\, \Delta t,
\end{equation}
with a per-cell flag cleared after the first crossing so subsequent
re-crossings are ignored.

\paragraph{Early termination on P8.}
When \code{stopAfterPointActivation = true}, the cell nearest
\code{stopActivationPoint} (P8 in the Niederer benchmark) is monitored
and \code{effectiveEndTime} is set to $t_{\text{act}} +
$~\code{stopDelayAfterActivation} once its $t_{\text{act}} > 0$.

\paragraph{Diagonal sampling.}
\code{sampleIDW} performs inverse-distance-weighted interpolation at
\code{nDiagonalSamples} points between \code{diagonalStart} and
\code{diagonalEnd}, with $k = $~\code{sampleNeighbours},
$p = $~\code{samplePower}. Because a diagonal sample averages its $k$ nearest
cells and unactivated cells contribute $t_{\text{act}}=0$, a sample watched
\emph{live} climbs from below as the front fills its neighbourhood; the final
value (all $k$ neighbours activated) is the correct one.

\paragraph{Live activation output.}
\code{writeActivationSamples} is called once before the time loop (so
\code{diagonalActivationTime.dat} and \code{P8\_activationTime.dat} exist from
$t=0$ with every point pending) and again whenever the activated-cell count
increases, overwriting the files. The run can therefore be monitored as the
front advances; the driver exposes the live file in each run directory by a
symbolic link that is replaced by a real copy at the end.

\paragraph{Activation-failure detection.}
When \dictkey{abortOnActivationFailure = true} the solver stops early
(cleanly, exit~0) if the activation fails, avoiding a full run to
\code{endTime}. Two criteria, using the running activated-cell count $N_a$ and
the time $t_\star$ of the last new activation:
\begin{itemize}[topsep=2pt,itemsep=1pt]
  \item \textbf{no ignition:} $N_a=0$ and
        $t>\dictkey{activationIgnitionTimeout}$;
  \item \textbf{conduction block:} $N_a>0$, the stop point is not yet reached,
        and $t-t_\star>\dictkey{activationStallTime}$.
\end{itemize}
The outcome is always recorded in
\code{postProcessing/.../activationStatus.dat} as one of \code{OK},
\code{FAILED\_NO\_IGNITION}, \code{FAILED\_CONDUCTION\_BLOCK} or
\code{INCOMPLETE} (with the final time and $N_a$). The driver reads this file
and, on a failure, logs it to \code{activation\_status\_summary.dat} and
continues with the next configuration rather than aborting the sweep.

% ================================================================== %
\section{Solver flow chart}
% ================================================================== %

The pseudo-code below follows the actual solver.  The shared setup is
common to all nonlinear methods; then one of the three method-specific
blocks is executed according to \dictkey{nonlinearMethod}.

\begin{algorithm}[H]
\caption{Shared time-step setup and final commit/rollback.}
\begin{algorithmic}[1]
\State Read dictionaries and assemble $\Mmass$, $\Kstif$, $\Lop=\Kstif/(\chi C_m)$.
\While{$t^n < T$}
  \State Save snapshots: $\mathsf{V}^n$, $\Iion^n$, $\mathbf{s}^n$.
  \State Build $\Aimp=\Mmass/\Delta t-\theta\Lop$ and
         $\Bimp=\Mmass/\Delta t+(1-\theta)\Lop$.
  \State Build old source $S^n=-\Iion^n+\Iext^n/(\chi C_m)$.
  \If{\dictkey{jfnkInitGuessOrder}$\ge 1$ and previous step exists}
    \State $\mathsf{x}_0\gets\mathsf{V}^n+
      (\Delta t/\Delta t_{\rm prev})(\mathsf{V}^n-\mathsf{V}^{n-1})$.
    \State Clamp $\mathsf{x}_0$ to
      $[\dictkey{jfnkInitGuessVmMin},\dictkey{jfnkInitGuessVmMax}]$.
  \Else
    \State $\mathsf{x}_0\gets\mathsf{V}^n$.
  \EndIf
  \State Run Algorithm~\ref{alg:picard},~\ref{alg:diagonal} or~\ref{alg:jfnk}.
  \If{not converged and \dictkey{nonlinearAcceptUnconverged} is false}
    \State Roll back $\Vm,\Iion,\mathbf{s}$ to the saved $t^n$ snapshots.
  \Else
    \State Commit the final nonlinear iterate.
  \EndIf
  \State Update activation times and benchmark samples.
\EndWhile
\end{algorithmic}
\end{algorithm}

\begin{algorithm}[H]
\caption{Picard nonlinear solve.}
\label{alg:picard}
\begin{algorithmic}[1]
\Require Initial candidate $\mathsf{x}_0$, state snapshot $\mathbf{s}^n$.
\For{$k=0,\dots,\code{nonlinearIterations}-1$}
  \State Store previous iterate $\mathsf{x}_{k-1}$, $I_{\mathrm{ion},k-1}$,
         $\mathbf{s}_{k-1}$ for residuals.
  \State \Call{evaluateNonlinearFields}{$\mathsf{x}_k$}:
    reset states to $\mathbf{s}^n$, reconstruct $\Vm$ at Iion quadrature
    points if enabled, call \code{ionicModel->solveODE}, compute $I_{\mathrm{ion},k}$.
  \State Build RHS $\mathsf{r}(I_{\mathrm{ion},k})$ using~\eqref{eq:rhsIion}.
  \State Solve $\Aimp\,\mathsf{V}_*=\mathsf{r}(I_{\mathrm{ion},k})$ with
    \dictkey{linearSolver}.
  \State Relax:
    $\mathsf{x}_{k+1}\gets\mathsf{x}_{k}+
    \code{nonlinearRelaxation}(\mathsf{V}_*-\mathsf{x}_{k})$.
  \State Re-evaluate ODEs and $\Iion$ at $\mathsf{x}_{k+1}$.
  \State Compute $R_{\mathrm{cpl}}$, $R_{\Vm}$, $R_{\Iion}$ and optional
    $R_{\mathbf{s}}$.
  \If{all requested criteria pass and $k+1\ge\code{minNonlinearIterations}$}
    \State Mark converged and exit loop.
  \EndIf
\EndFor
\end{algorithmic}
\end{algorithm}

\begin{algorithm}[H]
\caption{\code{diagonalIion} nonlinear solve.}
\label{alg:diagonal}
\begin{algorithmic}[1]
\Require Initial candidate $\mathsf{x}_0$, state snapshot $\mathbf{s}^n$.
\For{$k=0,\dots,\code{nonlinearIterations}-1$}
  \State Evaluate states and $I_{\mathrm{ion},k}$ at $\mathsf{x}_k$ exactly as in Picard.
  \State Approximate the local derivative
    $S'_k\approx-[\Iion(\mathsf{x}_k+\varepsilon_d)-\Iion(\mathsf{x}_k)]/\varepsilon_d$,
    with $\varepsilon_d=\dictkey{diagonalIionEpsilon}$.
  \State Assemble
    $\Aimp_k=\Aimp-\theta\,\mathrm{diag}(S'_k)$.
  \State Build corrected RHS
    $\mathsf{r}_k=\mathsf{r}(I_{\mathrm{ion},k})-\theta\,\mathrm{diag}(S'_k)\mathsf{x}_k$.
  \State Solve $\Aimp_k\mathsf{V}_*=\mathsf{r}_k$ with \dictkey{linearSolver}.
  \State Relax, re-evaluate ODEs/$\Iion$, compute residuals and test convergence
    as in Picard.
\EndFor
\end{algorithmic}
\end{algorithm}

\begin{algorithm}[H]
\caption{Jacobian-Free Newton--Krylov solve.}
\label{alg:jfnk}
\begin{algorithmic}[1]
\Require Initial candidate $\mathsf{x}_0$, state snapshot $\mathbf{s}^n$.
\For{$k=0,\dots,\code{nonlinearIterations}-1$}
  \State Evaluate $\mathsf{F}_k=\mathsf{F}(\mathsf{x}_k)$; this resets
    states to $\mathbf{s}^n$, solves ODEs and recomputes $\Iion$.
  \State Compute residuals. If all requested criteria pass, mark converged
    and exit.
  \State Define matrix-free product
    $\mathcal{J}_k\mathbf{v}=
    [\mathsf{F}(\mathsf{x}_k+\varepsilon\mathbf{v})-\mathsf{F}_k]/\varepsilon$,
    with $\varepsilon$ from \dictkey{jfnkEpsilon}.
  \State Solve $\mathcal{J}_k\boldsymbol{\delta}_k=-\mathsf{F}_k$ by
    restarted GMRES using \dictkey{jfnkMaxKrylovIterations},
    \dictkey{jfnkMaxRestarts} and \dictkey{jfnkLinearTolerance}.
  \If{\dictkey{jfnkLineSearch}}
    \State Start $\alpha=\code{nonlinearRelaxation}$ and backtrack until
      Armijo~\eqref{eq:armijo} succeeds, the trial budget is exhausted,
      or \dictkey{jfnkLineSearchAlphaMin} is reached.
  \Else
    \State $\alpha\gets\code{nonlinearRelaxation}$.
  \EndIf
  \State $\mathsf{x}_{k+1}\gets\mathsf{x}_{k}+\alpha\boldsymbol{\delta}_k$.
  \State Re-evaluate $\mathsf{F}(\mathsf{x}_{k+1})$, update fields, compute
    residuals and test convergence.
\EndFor
\end{algorithmic}
\end{algorithm}

\noindent\textbf{Legend.} \pyconst{Blue} = Python constant set in
\code{run\_convergence.py}; \dictkey{green} = OpenFOAM dictionary key
read by the solver. Constants and keys are mapped one-to-one in
section~\ref{sec:driver}.

% ================================================================== %
\section{Manufactured driver: \texttt{run\_convergence.py}}
\label{sec:driver}
% ================================================================== %

The convergence script
\code{/home/pablo/OpenFOAM/pablo-v2312/run/tutorials\_electro/highOrderManufacturedFDAImplicitJfNK/run\_convergence.py}
generates an MMS case whose numerical path is intended to match this
solver. The tables below describe the constants exposed there, the
OpenFOAM dictionary entries they generate, and which part of the
discrete equations they affect. The same dictionary keys are read by
\code{highOrderElectroActivationFoamImplicitJfNK}; only problem-specific
items such as exact fields versus TNNP/stimulus differ.

\subsection{Geometry and parameter sweep}
\begin{longtable}{@{}p{0.30\linewidth}p{0.28\linewidth}p{0.36\linewidth}@{}}
\toprule
Python constant & OpenFOAM key & Affected equation / mechanism \\
\midrule \endhead
\code{DIMENSIONS} & --- &
   Selects 2D or 3D manufactured domain and the corresponding LRE
   stencils in \code{P\_LEVELS\_BY\_DIM}. \\
\code{N\_VALUES} & \code{blockMeshDict} cell counts &
   Spatial refinement. This changes FV/LRE discretisation error in the
   PDE, Iion quadrature and state reconstruction. \\
\code{DT\_VALUES} & \code{controlDict.deltaT} &
   Outer PDE time step. Also sets the integration interval for every
   local ODE solve. \\
\code{END\_TIME} & \code{controlDict.endTime} &
   Final time of the transient MMS test. \\
\code{IMPLICIT\_SCHEMES} &
   \code{spatialIntegrationProperties.implicitScheme} &
   Chooses $\theta$ in~\eqref{eq:theta}: \code{backwardEuler}
   ($\theta=1$) or \code{crankNicolson} ($\theta=1/2$). \\
\code{MASS\_MATRIX\_VALUES} &
   \code{spatialIntegrationProperties.massMatrix} &
   Chooses lumped or LRE-consistent $\Mmass$ in
   $\Aimp=\Mmass/\Delta t-\theta\Lop$. \\
\code{ACTIVE\_HO\_PAIRS\_Vm\_States} &
   \code{useHighOrder\_Vm} / \code{useHighOrder\_Iion}
   + \code{LRECoeffs\_*} &
   Tuple \code{(Vm-level, Iion/states-level)}. Vm level affects
   diffusion and consistent mass; Iion/states level affects where ODEs
   are solved and where the nonlinear source is integrated. \\
\code{P\_LEVELS\_BY\_DIM} & \code{highOrderCoeffs} &
   Maps \code{p1,p2,p3} to LRE polynomial order \code{N}, stencil
   target \code{Nn}, radial weight exponent \code{k}, and
   \code{maxStencilSize}. \\
\bottomrule
\end{longtable}

\subsection{Nonlinear method and convergence controls}
\begin{longtable}{@{}p{0.35\linewidth}p{0.30\linewidth}p{0.30\linewidth}@{}}
\toprule
Python constant & OpenFOAM key & Affected equation / mechanism \\
\midrule \endhead
\code{NONLINEAR\_METHOD} & \code{nonlinearMethod} &
   Selects the nonlinear strategy for the PDE/source coupling:
   \code{Picard}, \code{JFNK}, or \code{diagonalIion}. \\
\code{IMPLICIT\_NONLINEAR\_ITERATIONS} & \code{nonlinearIterations} &
   Maximum number of nonlinear corrections per PDE time step. \\
\code{IMPLICIT\_MIN\_NONLINEAR\_ITERATIONS} &
   \code{minNonlinearIterations} &
   Minimum number of corrections before a convergence test may stop the
   loop. \\
\code{NONLINEAR\_TOLERANCE} & \code{nonlinearTolerance} &
   Coupled PDE residual tolerance $R_{\mathrm{cpl}}$. Used directly by
   JFNK and diagonalIion; in the electro solver it is also reported for
   Picard. \\
\code{NONLINEAR\_VM\_TOLERANCE} & \code{nonlinearVmTolerance} &
   Relative correction of the PDE unknown $\Vm$. \\
\code{NONLINEAR\_STATES\_TOLERANCE} &
   \code{nonlinearStatesTolerance} &
   Relative correction of all ODE states. In the MMS solver these are
   $u_1,u_2,u_3$; in this electro solver these are the TNNP state
   variables. \\
\code{NONLINEAR\_REQUIRE\_STATES\_CONVERGENCE} &
   \code{nonlinearRequireStatesConvergence} &
   If true, the state residual is part of the stopping criterion. \\
\code{NONLINEAR\_ACCEPT\_UNCONVERGED} &
   \code{nonlinearAcceptUnconverged} &
   If false, a failed electro nonlinear solve rolls back $\Vm,\Iion$ and
   states to $t^n$. \\
\code{NONLINEAR\_IION\_TOLERANCE} & \code{nonlinearIionTolerance} &
   Relative correction of the nonlinear source $\Iion$. \\
\code{NONLINEAR\_RELAXATION} & \code{nonlinearRelaxation} &
   Picard/diagonal mixing factor or initial Newton step length for
   JFNK/Armijo. \\
\bottomrule
\end{longtable}

\subsection{JFNK and diagonal-Iion controls}
\begin{longtable}{@{}p{0.35\linewidth}p{0.30\linewidth}p{0.30\linewidth}@{}}
\toprule
Python constant & OpenFOAM key & Affected equation / mechanism \\
\midrule \endhead
\code{JFNK\_MAX\_KRYLOV\_ITERATIONS} & \code{jfnkMaxKrylovIterations} &
   Restart length $m$ of GMRES(m), controlling memory and Krylov
   approximation quality. \\
\code{JFNK\_MAX\_RESTARTS} & \code{jfnkMaxRestarts} &
   Number of times the GMRES basis may be discarded and rebuilt. \\
\code{JFNK\_LINEAR\_TOLERANCE} & \code{jfnkLinearTolerance} &
   Relative tolerance of the Newton correction solve
   $\mathsf{J}\delta=-\mathsf{F}$. \\
\code{JFNK\_EPSILON} & \code{jfnkEpsilon} &
   Base perturbation for matrix-free Jacobian-vector products. It should
   be above the ODE/noise floor but small enough to approximate a
   derivative. \\
\code{JFNK\_INIT\_GUESS\_ORDER} & \code{jfnkInitGuessOrder} &
   Chooses lagged or linearly extrapolated initial $\Vm$ for the
   nonlinear solve. \\
\code{JFNK\_INIT\_GUESS\_VM\_MIN/MAX} &
   \code{jfnkInitGuessVmMin/Max} &
   Clamp interval for extrapolated guess and, if enabled, ODE inputs. \\
\code{JFNK\_CLAMP\_ODE\_INPUT} & \code{jfnkClampODEInput} &
   Protects the ODE solve from unphysical JFNK finite-difference probes. \\
\code{JFNK\_LINE\_SEARCH} & \code{jfnkLineSearch} &
   Turns Armijo backtracking on/off. \\
\code{JFNK\_LINE\_SEARCH\_MAX\_ITER} &
   \code{jfnkLineSearchMaxIter} &
   Maximum number of halvings of the Newton step. \\
\code{JFNK\_LINE\_SEARCH\_ALPHA\_MIN} &
   \code{jfnkLineSearchAlphaMin} &
   Floor for the Armijo step length. \\
\code{JFNK\_ARMIJO\_C} & \code{jfnkArmijoC} &
   Sufficient-decrease parameter in~\eqref{eq:armijo}. \\
\code{DIAGONAL\_IION\_EPSILON} & \code{diagonalIionEpsilon} &
   Finite-difference step for the diagonal source derivative in
   \code{diagonalIion}. \\
\bottomrule
\end{longtable}

\subsection{Linear PDE solver}
\begin{longtable}{@{}p{0.30\linewidth}p{0.30\linewidth}p{0.35\linewidth}@{}}
\toprule
Python constant & OpenFOAM key & Affected equation / mechanism \\
\midrule \endhead
\code{IMPLICIT\_LINEAR\_SOLVER} & \code{linearSolver} &
   Eigen sparse solver used by Picard and diagonalIion for
   $\Aimp\mathsf{V}=\mathsf{r}$. JFNK instead uses GMRES on the
   Jacobian-free shell. \\
\code{IMPLICIT\_TOLERANCE} & \code{tolerance} &
   BiCGSTAB tolerance for the linear PDE solve. \\
\code{IMPLICIT\_MAX\_ITERATIONS} & \code{maxIterations} &
   BiCGSTAB iteration cap. \\
\bottomrule
\end{longtable}

\subsection{State ODE integrator}
\begin{longtable}{@{}p{0.30\linewidth}p{0.30\linewidth}p{0.35\linewidth}@{}}
\toprule
Python constant & OpenFOAM key & Affected equation / mechanism \\
\midrule \endhead
\code{STATE\_ODE\_SOLVER} & \code{stateODESolver} /
   \code{timeIntegrationProperties.solver} &
   Local ODE integrator for MMS states or TNNP states. \\
\code{STATE\_ODE\_ABS\_TOL} & \code{stateODEAbsTol} /
   \code{timeIntegrationProperties.absTol} &
   Absolute tolerance of the ODE solve. \\
\code{STATE\_ODE\_REL\_TOL} & \code{stateODERelTol} /
   \code{timeIntegrationProperties.relTol} &
   Relative ODE tolerance. This controls the noise floor seen by JFNK
   finite differences. \\
\code{STATE\_ODE\_INITIAL\_STEP} & \code{stateODEInitialStep} /
   \code{initialODEStep} &
   First RKF45 sub-step attempted inside one PDE time step. \\
\code{STATE\_ODE\_MAX\_STEPS} & \code{stateODEMaxSteps} /
   \code{maxSteps} &
   Maximum ODE substeps per local solve. \\
\bottomrule
\end{longtable}

\subsection{LRE high-order parameters per p-level}
\code{P\_LEVELS\_BY\_DIM[dim][p]} gives, for each (dim, p):
\code{N} (polynomial order), \code{Nn} (stencil neighbours),
\code{k} (search radius), \code{maxStencilSize} (cap). The script
writes these into a nested
\code{highOrderCoeffs\{LRECoeffs\_Vm\{\dots\}; LRECoeffs\_Iion\{\dots\}\}}
block whenever HO is enabled for either Vm or $\Iion$.

\subsection{Diagnostics and run-control}
\begin{longtable}{@{}p{0.35\linewidth}p{0.30\linewidth}p{0.30\linewidth}@{}}
\toprule
Python constant & OpenFOAM key & Affected mechanism \\
\midrule \endhead
\code{WRITE\_NONLINEAR\_RESIDUALS} &
   \code{writeNonlinearResiduals} & Residual log on/off. \\
\code{KEEP\_LOGS\_ON\_SUCCESS} & --- &
   Python-only switch: keep solver/blockMesh logs after successful runs.
   Failed-run logs are always kept. \\
\bottomrule
\end{longtable}

For the electro-activation tutorial the analogous driver also exposes
activation-threshold, P8 stopping and diagonal-sampling parameters. They
do not affect the PDE/ODE solve itself; they affect only stopping and
post-processing of activation times.

% ================================================================== %
\section{Results}
\label{sec:results}
% ================================================================== %

This section collects the current convergence results for the 2D and
3D Niederer-style benchmarks, run with the three nonlinear strategies
exposed by the solver (\code{Picard}, \code{diagonalIion}, \code{JFNK}).
All runs use the Crank--Nicolson time scheme ($\theta = 1/2$), the
consistent high-order mass matrix when LRE is enabled, the
\code{RKF45} state ODE integrator, and a single requested
$\Delta t_{\text{req}} = 10^{-4}$~s. Six LRE pairs $(\Vm,\Iion)$ are
sampled: \code{NO}/\code{NO}, \code{NO}/\code{p1}, \code{NO}/\code{p3},
\code{p3}/\code{NO}, \code{p3}/\code{p1}, \code{p3}/\code{p3}. Each
table lists, for the metric in its caption, the value of every
available run; \textsc{NC} marks combinations that have not been run
yet. The 2D results live under
\path{highOrderNiedererEtAl2012ImplicitJfNK[_diagonalIion|_Picard]/results/example0/2D/};
the 3D results under the same trees but with the \path{3D/} suffix.
Wall-clock time is the elapsed time reported by
\texttt{/usr/bin/time -v}, rounded to integer minutes. Peak RSS is
the maximum resident set size of the solver process, expressed in
MiB ($\lfloor\text{kB}/1024\rfloor$).

\subsection{2D Niederer-style benchmark}
\label{sec:results:2D}

The 2D case follows the 20$\times$8~mm planar reduction used by the
explicit and non-JfNK implicit reports (single cell in $z$, planar
fibres), exercised at mesh sizes
$\Delta x\in\{0.5, 0.25, 0.125, 0.1\}$~mm; only the \code{Picard} sweep
currently includes the finest mesh.

\begin{center}\footnotesize
\begin{longtable}{llrrrr}
\caption{2D P8 activation time [ms] at $\Delta t = 10^{-4}$~s. Columns
are mesh sizes $\Delta x$ in mm.}\\
\toprule
Vm & Iion & $0.5$ & $0.25$ & $0.125$ & $0.1$ \\
\midrule\endfirsthead\toprule
Vm & Iion & $0.5$ & $0.25$ & $0.125$ & $0.1$ \\
\midrule\endhead
\multicolumn{6}{l}{\textit{\code{Picard}}}\\
NO & NO & 153.32 & 62.99 & 45.03 & 42.76 \\
NO & p1 &  91.05 & 54.86 & 43.21 & 41.84 \\
NO & p3 &  32.46 & 39.14 & 39.83 & NC    \\
p3 & NO & 180.55 & 68.84 & 48.83 & NC    \\
p3 & p1 &  68.92 & 58.81 & 46.33 & NC    \\
p3 & p3 &  33.09 & 40.22 & 40.84 & NC    \\\hline
\multicolumn{6}{l}{\textit{\code{diagonalIion}}}\\
NO & NO & 153.16 & 62.86 & 44.92 & NC    \\
NO & p1 &  90.92 & 54.74 & 43.08 & NC    \\
NO & p3 &  32.37 & 38.97 & 39.66 & NC    \\
p3 & NO & 180.42 & 68.71 & 48.67 & NC    \\
p3 & p1 &  68.83 & 58.70 & 46.20 & NC    \\
p3 & p3 &  33.00 & 40.03 & 40.65 & NC    \\\hline
\multicolumn{6}{l}{\textit{\code{JFNK}}}\\
NO & NO & 153.13 & 62.85 & 44.87 & NC    \\
NO & p1 &  90.92 & 54.72 & NC    & NC    \\
NO & p3 &  32.36 & 38.92 & 39.59 & NC    \\
p3 & NO & 180.38 & 68.69 & NC    & NC    \\
p3 & p1 &  68.83 & NC    & NC    & NC    \\
p3 & p3 &  33.00 & NC    & NC    & NC    \\
\bottomrule\end{longtable}\end{center}

\begin{center}\footnotesize
\begin{longtable}{llrrrr}
\caption{2D solver wall-clock time [min] at $\Delta t = 10^{-4}$~s
(rounded to integer minutes).}\\
\toprule
Vm & Iion & $0.5$ & $0.25$ & $0.125$ & $0.1$ \\
\midrule\endfirsthead\toprule
Vm & Iion & $0.5$ & $0.25$ & $0.125$ & $0.1$ \\
\midrule\endhead
\multicolumn{6}{l}{\textit{\code{Picard}}}\\
NO & NO &  15 &  31 &  83 & 125 \\
NO & p1 &  20 &  51 & 141 & 197 \\
NO & p3 &  23 & 102 & 402 & NC  \\
p3 & NO &  18 &  40 & 112 & NC  \\
p3 & p1 &  20 &  60 & 176 & NC  \\
p3 & p3 &  22 & 105 & 419 & NC  \\\hline
\multicolumn{6}{l}{\textit{\code{diagonalIion}}}\\
NO & NO &  18 &  31 &  81 & NC  \\
NO & p1 &  22 &  52 & 147 & NC  \\
NO & p3 &  29 & 131 & 533 & NC  \\
p3 & NO &  21 &  38 & 107 & NC  \\
p3 & p1 &  23 &  59 & 164 & NC  \\
p3 & p3 &  28 & 133 & 549 & NC  \\\hline
\multicolumn{6}{l}{\textit{\code{JFNK}}}\\
NO & NO &  59 & 155 & 544 & NC  \\
NO & p1 &  82 & 263 & NC  & NC  \\
NO & p3 & 124 & 647 & 3407 & NC \\
p3 & NO &  67 & 169 & NC  & NC  \\
p3 & p1 &  87 & NC  & NC  & NC  \\
p3 & p3 & 131 & NC  & NC  & NC  \\
\bottomrule\end{longtable}\end{center}

\begin{center}\footnotesize
\begin{longtable}{llrrrr}
\caption{2D peak resident set size [MiB] at $\Delta t = 10^{-4}$~s.}\\
\toprule
Vm & Iion & $0.5$ & $0.25$ & $0.125$ & $0.1$ \\
\midrule\endfirsthead\toprule
Vm & Iion & $0.5$ & $0.25$ & $0.125$ & $0.1$ \\
\midrule\endhead
\multicolumn{6}{l}{\textit{\code{Picard}}}\\
NO & NO & 148 & 152 & 174 & 194 \\
NO & p1 & 149 & 157 & 191 & 219 \\
NO & p3 & 155 & 177 & 288 & NC  \\
p3 & NO & 155 & 191 & 346 & NC  \\
p3 & p1 & 156 & 194 & 358 & NC  \\
p3 & p3 & 161 & 218 & 457 & NC  \\\hline
\multicolumn{6}{l}{\textit{\code{diagonalIion}}}\\
NO & NO & 148 & 153 & 176 & NC  \\
NO & p1 & 149 & 158 & 194 & NC  \\
NO & p3 & 155 & 181 & 296 & NC  \\
p3 & NO & 155 & 191 & 347 & NC  \\
p3 & p1 & 156 & 195 & 365 & NC  \\
p3 & p3 & 161 & 218 & 457 & NC  \\\hline
\multicolumn{6}{l}{\textit{\code{JFNK}}}\\
NO & NO & 148 & 153 & 173 & NC  \\
NO & p1 & 149 & 158 & NC  & NC  \\
NO & p3 & 155 & 181 & 299 & NC  \\
p3 & NO & 152 & 187 & NC  & NC  \\
p3 & p1 & 153 & NC  & NC  & NC  \\
p3 & p3 & 160 & NC  & NC  & NC  \\
\bottomrule\end{longtable}\end{center}

\paragraph{Discussion of the 2D results.}
At fixed $(\Delta x, \Delta t, $LRE pair$)$ the three nonlinear
strategies give activation times that agree to within $\sim 0.1$--$0.2$~ms;
e.g.\ at $\Delta x = 0.125$~mm with $(\code{NO}, \code{NO})$ the values
are $45.03/44.92/44.87$~ms for Picard, diagonalIion, JFNK respectively.
The 2D nonlinear-coupling choice is therefore essentially \emph{free}
from the accuracy standpoint here --- the relevant errors are spatial
and temporal, not nonlinear-iteration errors. The three methods differ
sharply in cost: JFNK is $\sim 3$--$6\times$ more expensive than
Picard for the same case at $\Delta x = 0.5$~mm and grows faster than
linearly with mesh size, because every Krylov matrix--vector product
performs a full ODE re-evaluation
($\Delta x = 0.125$, $(\code{NO},\code{p3})$: $402$~min for Picard,
$533$~min for diagonalIion, $3407$~min --- $\sim 2.4$~days --- for
JFNK). Peak RSS is dominated by the LRE stencil and is essentially
method-independent; the JFNK column is slightly lower for some entries
because the matrix-free shell avoids the diagonal correction and the
explicit BiCGSTAB workspace.

\subsection{3D Niederer-style benchmark}
\label{sec:results:3D}

The 3D sweep follows the published Niederer slab geometry
($20\times 3\times 7$~mm, $40\times 6\times 14$ cells at
$\Delta x = 0.5$~mm). At the time of writing only the
\code{diagonalIion} and \code{Picard} sweeps include 3D runs, all at
$\Delta x = 0.5$~mm. The \code{JFNK} 3D sweep has not yet been run.

\begin{center}\footnotesize
\begin{longtable}{llrr}
\caption{3D P8 activation time [ms] at $\Delta x = 0.5$~mm,
$\Delta t = 10^{-4}$~s. Columns are the available nonlinear methods.}\\
\toprule
Vm & Iion & \code{Picard} & \code{diagonalIion} \\
\midrule\endfirsthead\toprule
Vm & Iion & \code{Picard} & \code{diagonalIion} \\
\midrule\endhead
NO & NO & 136.67 & 136.61 \\
NO & p1 &  33.50 &  33.39 \\
NO & p3 &  30.59 &  30.51 \\
p3 & NO & 149.13 & 149.06 \\
p3 & p1 &  32.44 &  32.36 \\
p3 & p3 &  29.63 & NC     \\
\bottomrule\end{longtable}\end{center}

\begin{center}\footnotesize
\begin{longtable}{llrr}
\caption{3D solver wall-clock time [min] at $\Delta x = 0.5$~mm,
$\Delta t = 10^{-4}$~s.}\\
\toprule
Vm & Iion & \code{Picard} & \code{diagonalIion} \\
\midrule\endfirsthead\toprule
Vm & Iion & \code{Picard} & \code{diagonalIion} \\
\midrule\endhead
NO & NO &   80 &   82 \\
NO & p1 &  168 &  229 \\
NO & p3 & 1489 & 2159 \\
p3 & NO &  163 &  143 \\
p3 & p1 &  167 &  269 \\
p3 & p3 & 1487 & NC   \\
\bottomrule\end{longtable}\end{center}

\begin{center}\footnotesize
\begin{longtable}{llrr}
\caption{3D peak resident set size [MiB] at $\Delta x = 0.5$~mm,
$\Delta t = 10^{-4}$~s.}\\
\toprule
Vm & Iion & \code{Picard} & \code{diagonalIion} \\
\midrule\endfirsthead\toprule
Vm & Iion & \code{Picard} & \code{diagonalIion} \\
\midrule\endhead
NO & NO &  156 &  157 \\
NO & p1 &  211 &  211 \\
NO & p3 &  787 &  849 \\
p3 & NO & 1123 & 1121 \\
p3 & p1 & 1164 & 1162 \\
p3 & p3 & 1565 & NC   \\
\bottomrule\end{longtable}\end{center}

\paragraph{Discussion of the 3D results.}
On the coarse 3D slab ($\Delta x = 0.5$~mm), activating LRE on the
ionic source alone changes the activation time from $\approx 136.7$~ms
(\code{NO}/\code{NO}) to $\approx 33.5$~ms (\code{NO}/\code{p1}) and
$\approx 30.6$~ms (\code{NO}/\code{p3}) --- the same qualitative
high-order-on-source effect already seen in the 2D sweep and in the
explicit/implicit Niederer reports. The fully high-order
(\code{p3}/\code{p3}) reaches $29.63$~ms with Picard (the diagonalIion
counterpart has not yet been completed). The two nonlinear methods
again agree to better than $0.1$~ms for every entry where both are
available; Picard is generally faster, with the most striking gap on
the heaviest case (\code{NO}/\code{p3}: $1489$~min for Picard vs.\
$2159$~min for diagonalIion). Peak RSS scales with the LRE stencil
size and is essentially method-independent. The dataset is too thin
in $\Delta x$ to draw quantitative conclusions about mesh convergence
in 3D; the immediate gap to fill is the missing \code{p3}/\code{p3}
diagonalIion entry and a JFNK 3D sweep at the same mesh.

% ================================================================== %
\appendix
\section{Dimension table}
% ================================================================== %

\begin{center}
\begin{tabular}{@{}lll@{}}
\toprule
Field & OpenFOAM \code{dimensionSet} & SI \\
\midrule
\code{Vm} & $[1, 2, -3, 0, 0, -1, 0]$ & volt \\
\code{Iion} & \code{dimVoltage/dimTime} & \si{\volt\per\second} \\
\code{externalStimulusCurrent} & \code{dimCurrent/dimVolume} &
    \si{\ampere\per\meter\cubed} \\
\code{conductivity} &
    \code{dimCurrent\^{}2 dimTime\^{}3 / (dimMass dimVolume)} &
    \si{\siemens\per\meter} \\
\code{chi} & \code{dimless/dimLength} & \si{\per\meter} \\
\code{Cm} & \code{dimCurrent dimTime / (dimVoltage dimArea)} &
    \si{\farad\per\meter\squared} \\
\bottomrule
\end{tabular}
\end{center}

\section{File map}

\begin{center}
\begin{tabular}{@{}p{0.55\linewidth}p{0.40\linewidth}@{}}
\toprule
Path (relative to \code{cardiacFoam-pc/}) & Role \\
\midrule
\code{applications/solvers/.../JfNK.C} & Main translation unit, including
   the embedded TNNP CellML kernel, state update, current evaluation, and
   local ODE integration. \\
\code{applications/solvers/.../createFields.H} & Field initialisation
   and dictionary reading. \\
\code{tutorials\_electro/.../run\_convergence.py} & Driver script. \\
\code{constant/cardiacProperties} & $\chi$, $C_{m}$, $\Dten$,
   ionic model. \\
\code{constant/timeIntegrationProperties} & ODE solver, tolerances,
   sub-step cap. \\
\code{constant/spatialIntegrationProperties} & Time scheme, HO flags,
   JFNK / linear / line-search options. \\
\code{constant/stimulusProtocol} & Stimulus regions and timing. \\
\bottomrule
\end{tabular}
\end{center}

% ================================================================== %
\section{Differences vs.\ the PETSc sibling}
% ================================================================== %

Algorithmically the two solvers are identical. The only differences
are in the linear-algebra back-end and a few user-visible knobs:

\begin{center}
\begin{tabular}{@{}p{0.42\linewidth}p{0.26\linewidth}p{0.26\linewidth}@{}}
\toprule
Aspect & This solver (Eigen) & PETSc sibling \\
\midrule
Sparse matrix container & \code{Eigen::SparseMatrix} & PETSc \code{Mat} \\
Sparse direct solver & \code{Eigen::SparseLU} & \code{KSPPREONLY+PCLU} \\
Iterative solver (Picard branch) & \code{Eigen::BiCGSTAB+ILUT} &
    \code{KSP+PC} (full PETSc catalogue) \\
JFNK inner Krylov & in-house GMRES & \code{KSPGMRES} on \code{MatShell} \\
Preconditioner inside JFNK & none (Krylov-only) & any PETSc PC \\
Build-time dependency & header-only Eigen3 &
    PETSc + MPI \\
Parallel run support & serial only & MPI-parallel out of the box \\
\bottomrule
\end{tabular}
\end{center}

All convergence-related parameters (\code{nonlinearTolerance},
\code{jfnkEpsilon}, \code{jfnkLineSearch*}, \code{jfnkInitGuess*},
\code{jfnkClampODEInput}, \code{nonlinearAcceptUnconverged},
\code{ODE\_*}) have the same names and meanings in both solvers, so
the same \code{run\_convergence.py} configuration produces an
algorithmically equivalent run with either back-end (modulo the
performance and parallel-scaling differences).

% ================================================================== %
\begin{thebibliography}{9}
% ================================================================== %

\bibitem{ttp2004} K.~H.~W.~J. ten Tusscher, D. Noble, P.~J. Noble, A.~V.
    Panfilov, \emph{A model for human ventricular tissue},
    AJP -- Heart and Circulatory Physiology
    \textbf{286} (2004) H1573--H1589.

\bibitem{niederer2012} S.~A. Niederer et al.,
    \emph{Verification of cardiac tissue electrophysiology simulators
    using an N-version benchmark}, Prog.\ Biophys.\ Mol.\ Biol.
    \textbf{104} (2012) 22--48.

\bibitem{knoll2004} D.~A. Knoll, D.~E. Keyes,
    \emph{Jacobian-Free Newton-Krylov methods: a survey of approaches
    and applications}, J.\ Comp.\ Phys.\ \textbf{193}
    (2004) 357--397.

\bibitem{noceWright} J. Nocedal, S.~J. Wright,
    \emph{Numerical Optimization}, 2nd edition, Springer, 2006.

\bibitem{saadbook} Y. Saad,
    \emph{Iterative Methods for Sparse Linear Systems}, 2nd edition,
    SIAM, 2003.

\bibitem{hairer2} E. Hairer, G. Wanner,
    \emph{Solving Ordinary Differential Equations II}, Springer, 1996.

\end{thebibliography}

\end{document}
